\chapterimage{figs/chapterhead.png}
\chapter{Download, instalação e compilação}
\label{chap:download_instalacao_compilacao}

\section{Introdução}
\label{sec:cap2_introducao}

Depois de compreender o papel do \textit{GenESyS} e a forma como este manual foi organizado, o leitor precisa colocá-lo em funcionamento. Esse é o objetivo deste capítulo. Em vez de tratar a compilação como tema central do livro, ela será tratada aqui como aquilo que realmente é para o usuário: uma etapa necessária para alcançar a aplicação gráfica e começar a trabalhar com modelos.

É importante notar que o \textit{GenESyS} é um projeto de software mais amplo do que uma única aplicação. Há kernel, parser, \textit{plugins}, testes, aplicações auxiliares e outros elementos internos. No entanto, essa complexidade não precisa ser absorvida pelo usuário logo no início. A meta prática desta etapa é mais simples: obter o projeto, preparar o ambiente, compilar a GUI e validar a primeira execução da aplicação gráfica. Se isso for alcançado com segurança, o leitor já terá o necessário para avançar para a fase de uso efetivo do simulador.

\section{Obtendo o projeto}
\label{sec:cap2_obtendo_projeto}

O \textit{GenESyS} está disponível como repositório de código-fonte. Assim, o primeiro passo consiste em obter uma cópia local do projeto. Em um ambiente com \texttt{git} devidamente instalado, esse procedimento é direto.

\begin{lstlisting}[language=bash,caption={Clonagem inicial do repositório do GenESyS},label={lst:cap2_git_clone}]
    git clone https://github.com/rlcancian/Genesys-Simulator.git
    cd Genesys-Simulator
\end{lstlisting}

Esse comando produz uma cópia local do repositório. A partir daí, o usuário passa a ter em sua máquina o conjunto de arquivos necessários para preparar o ambiente de trabalho. Convém, nesse ponto, fazer uma inspeção inicial da raiz do projeto, apenas para reconhecer sua presença e seu volume estrutural.

\begin{lstlisting}[language=bash,caption={Inspeção inicial da raiz do repositório},label={lst:cap2_inspecao_inicial}]
    pwd
    ls
\end{lstlisting}

O usuário não precisa ainda interpretar em profundidade tudo o que aparece nessa listagem. Basta reconhecer que há uma árvore-fonte, há uma pasta de modelos e há arquivos de configuração que sustentam a compilação e o restante do projeto.

\section{O que o usuário realmente precisa nessa fase}
\label{sec:cap2_o_que_usuario_precisa}

Uma dificuldade frequente em projetos acadêmicos e de pesquisa é supor que o usuário precise compreender toda a organização interna antes de conseguir utilizá-los. Para o \textit{GenESyS}, essa estratégia seria ruim. O usuário desta primeira parte do livro não precisa estudar imediatamente a estrutura do kernel, os \textit{plugins} ou o parser. Ele precisa apenas do que for suficiente para:
\begin{itemize}
    \item preparar um ambiente coerente;
    \item compilar a aplicação gráfica;
    \item iniciar a GUI;
    \item confirmar que a interface está funcional.
\end{itemize}

Portanto, a organização do repositório deve ser lida, neste momento, apenas com uma preocupação prática. A pasta \texttt{source/} contém o código da aplicação e do sistema. A pasta \texttt{models/} será importante nos próximos capítulos, pois é dela que virão os modelos de exemplo usados pelo usuário. Demais elementos do projeto podem permanecer, por ora, em segundo plano.

\section{Dependências do ambiente}
\label{sec:cap2_dependencias_ambiente}

A versão atual do \textit{GenESyS} utiliza \texttt{CMake} como sistema principal de configuração e build. Isso significa que o ambiente de compilação do usuário precisa oferecer um conjunto mínimo de ferramentas compatíveis com essa organização. Além do próprio \texttt{CMake}, o projeto depende de um compilador \texttt{C++} moderno e de uma instalação funcional de Qt para a aplicação gráfica.

Do ponto de vista prático, convém pensar o ambiente em três camadas:

\begin{enumerate}
    \item ferramentas de build;
    \item bibliotecas gráficas e de apoio;
    \item compilador e utilitários gerais de desenvolvimento.
\end{enumerate}

Na primeira camada, o elemento central é o \texttt{CMake}. Na segunda, o ponto mais importante é o Qt, especialmente os módulos associados à interface gráfica. Na terceira, entram o compilador \texttt{C++} e as ferramentas usuais de desenvolvimento da plataforma escolhida.

No caso específico da GUI, o projeto procura preferencialmente por \texttt{Qt6} e, se necessário, admite fallback para \texttt{Qt5}. Para o usuário, isso significa que o ambiente precisa disponibilizar ao menos uma dessas instalações de forma coerente.

\subsection{Checklist mínimo do ambiente}
\label{subsec:cap2_checklist_minimo_ambiente}

Antes de partir para a configuração do build, o usuário deve confirmar que possui, no mínimo:
\begin{itemize}
    \item \texttt{git};
    \item \texttt{cmake};
    \item um compilador \texttt{C++} moderno;
    \item uma instalação funcional de \texttt{Qt};
    \item ferramentas básicas de terminal e sistema de arquivos.
\end{itemize}

Esse checklist não substitui instruções específicas de cada plataforma, mas evita que o processo de compilação seja iniciado em ambiente claramente incompleto.

\section{Compilando a aplicação gráfica}
\label{sec:cap2_compilando_aplicacao_grafica}

A compilação da GUI deve ser conduzida em um diretório de build separado da árvore-fonte. Essa prática é recomendável por duas razões. A primeira é manter o repositório limpo, sem misturar arquivos gerados com os arquivos originais do projeto. A segunda é tornar mais fácil a repetição do processo em caso de ajustes, limpeza ou reconfiguração do ambiente.

\begin{lstlisting}[language=bash,caption={Criação do diretório de build},label={lst:cap2_criacao_build}]
    mkdir -p build
    cd build
\end{lstlisting}

Com o diretório de build criado, o usuário pode configurar o projeto. Como o caminho principal desta primeira parte do livro é a GUI, a configuração deve habilitar explicitamente a aplicação gráfica. Uma configuração inicial típica pode ser feita assim:

\begin{lstlisting}[language=bash,caption={Configuração inicial do build da GUI},label={lst:cap2_configuracao_gui}]
    cmake .. \
    -DGENESYS_BUILD_GUI_APPLICATION=ON
\end{lstlisting}

Dependendo da plataforma e do ambiente local, o usuário pode querer ajustar outras opções. No entanto, como orientação inicial do manual, o ponto essencial é garantir que a GUI seja incluída no processo de build. Depois disso, a compilação propriamente dita pode ser realizada.

\begin{lstlisting}[language=bash,caption={Compilação da GUI do GenESyS},label={lst:cap2_compilacao_gui}]
    cmake --build . --target genesys_qt_gui_application -j
\end{lstlisting}

Em alguns ambientes, o alvo agregado da GUI também pode ser usado, quando disponível:

\begin{lstlisting}[language=bash,caption={Alternativa de compilação pelo alvo agregado da GUI},label={lst:cap2_compilacao_gui_agregada}]
    cmake --build . --target genesys_gui -j
\end{lstlisting}

\subsection{Por que a GUI é explicitamente habilitada}
\label{subsec:cap2_por_que_gui_habilitada}

É importante compreender a lógica dessa configuração. O projeto contém mais de uma aplicação possível, mas o manual do usuário não será guiado por todas elas. O alvo desta primeira metade do livro é a interface gráfica. Por isso, a GUI é habilitada explicitamente na configuração do build. Essa decisão reduz ambiguidades e alinha o processo técnico à forma real como o usuário trabalhará com o sistema.

\section{Como interpretar a configuração e a compilação}
\label{sec:cap2_como_interpretar_configuracao_compilacao}

Em um primeiro contato, a saída textual do \texttt{CMake} e do processo de compilação pode parecer excessivamente detalhada. Contudo, o usuário não precisa interpretar tudo. Há apenas alguns sinais essenciais que devem ser observados.

Na configuração:
\begin{itemize}
    \item o Qt deve ser localizado corretamente;
    \item o processo não deve terminar com erro fatal;
    \item o diretório de build deve ser gerado normalmente.
\end{itemize}

Na compilação:
\begin{itemize}
    \item o processo deve completar a geração do executável da GUI;
    \item não devem aparecer falhas de linkedição ou de bibliotecas ausentes;
    \item o alvo da aplicação gráfica deve ser produzido ao final.
\end{itemize}

Essa leitura seletiva da saída do build é importante. O usuário não precisa ainda analisar toda a cadeia interna de alvos do projeto. Ele precisa apenas confirmar que o caminho até a aplicação gráfica foi bem-sucedido.

\section{O que fazer em caso de erro}
\label{sec:cap2_o_que_fazer_em_caso_de_erro}

Problemas de compilação fazem parte do trabalho com software, especialmente quando o projeto está em evolução. Isso não significa, porém, que o usuário precise entrar imediatamente em depuração arquitetural do sistema. A abordagem mais útil é começar pelas causas mais simples.

Em caso de erro, verifique primeiro:
\begin{itemize}
    \item se o Qt foi instalado corretamente;
    \item se o \texttt{cmake} está disponível e atualizado;
    \item se o compilador \texttt{C++} está corretamente configurado;
    \item se o diretório de build foi criado limpo;
    \item se não há resíduos de configurações anteriores incompatíveis.
\end{itemize}

Uma prática particularmente útil é remover o diretório \texttt{build} e repetir o processo em ambiente limpo quando surgirem inconsistências difíceis de interpretar. Em muitos casos, isso elimina resíduos de configuração ou escolhas anteriores inadequadas.

\section{Primeira execução da GUI}
\label{sec:cap2_primeira_execucao_gui}

Concluída a compilação, o próximo objetivo é iniciar a aplicação gráfica. O executável gerado corresponde ao alvo \texttt{genesys\_qt\_gui\_application}. Em um cenário típico, sua execução inicial será feita diretamente a partir do diretório de build ou do caminho exato em que o gerador utilizado o tiver produzido.

\begin{lstlisting}[language=bash,caption={Exemplo genérico de primeira execução da GUI},label={lst:cap2_primeira_execucao_gui}]
    ./genesys_qt_gui_application
\end{lstlisting}

O caminho exato pode variar, mas a intenção é sempre a mesma: chegar à janela principal da aplicação e confirmar que a GUI está funcional.

\subsection{O que validar na primeira execução}
\label{subsec:cap2_o_que_validar_primeira_execucao}

Ao iniciar a aplicação pela primeira vez, o usuário deve verificar alguns pontos muito concretos:

\begin{itemize}
    \item a janela principal abre corretamente;
    \item menus e áreas visíveis da interface aparecem de forma coerente;
    \item não há falha estrutural imediata ao iniciar a aplicação;
    \item a GUI parece pronta para abrir ou criar modelos.
\end{itemize}

Ainda não é necessário abrir modelos ou executar simulações. A meta desta etapa é mais modesta, mas fundamental: confirmar que o ambiente preparado realmente leva à aplicação gráfica em funcionamento.

\section{Primeiras observações úteis sobre a interface}
\label{sec:cap2_primeiras_observacoes_interface}

Uma vez que a GUI esteja aberta, o usuário já pode fazer uma leitura muito breve da aplicação, sem ainda tentar usá-la em profundidade. Convém observar:
\begin{itemize}
    \item a presença de barra de menus;
    \item a área central de trabalho;
    \item a existência de painéis, árvores ou abas laterais;
    \item a disponibilidade de comandos ligados a modelo, edição e simulação.
\end{itemize}

Essa observação não substitui o capítulo seguinte, que tratará especificamente da interface gráfica. Contudo, ela ajuda a transformar a primeira execução em experiência mais concreta do que simplesmente ``o programa abriu''.

\section{Preparando o próximo passo}
\label{sec:cap2_preparando_proximo_passo}

Se o leitor conseguiu chegar até aqui, então já possui o elemento mais importante desta fase inicial: uma instalação funcional do \textit{GenESyS} pronta para uso. Isso é suficiente para avançar. Não é necessário, neste momento, aprofundar-se nos detalhes da estrutura do código nem explorar todas as opções possíveis de build. A partir daqui, o manual passa do problema técnico da preparação do ambiente para o problema prático do uso do simulador.

Em outras palavras, o leitor já está em condição de começar a trabalhar com a aplicação gráfica propriamente dita.

\section{Considerações Finais}
\label{sec:cap2_consideracoes_finais}

O propósito deste capítulo foi conduzir o leitor desde a obtenção do projeto até a primeira execução funcional da GUI do \textit{GenESyS}. Embora essa etapa envolva download, dependências, configuração e compilação, seu papel dentro do livro continua sendo instrumental. O objetivo não era transformar o build no centro do manual, mas remover a barreira técnica necessária para que o uso do simulador pudesse começar de fato.

Com a aplicação gráfica em funcionamento, o próximo passo natural é compreender melhor a janela principal e sua lógica de organização. Por isso, o capítulo seguinte passa a tratar da interface gráfica como ambiente de trabalho do usuário, mostrando como a GUI organiza menus, áreas de edição, comandos de simulação, painéis de navegação e diferentes visões do modelo.

Teste seu entendimento desse conteúdo respondendo as seguintes perguntas:

\begin{enumerate}
    \item Por que a compilação da GUI é tratada neste manual como meio para alcançar o uso do simulador, e não como eixo organizador do livro?

    \item Que dependências mínimas o ambiente do usuário precisa fornecer antes da configuração do build?

    \item Por que é recomendável compilar o projeto em um diretório de build separado?

    \item Que sinais mostram que a primeira execução da aplicação gráfica foi bem-sucedida?
\end{enumerate}